Vision-based runway detection on approach is a promising, infrastructure-independent
building block for assisted and autonomous landing, but training data that covers
diverse airports, viewpoints and, above all, adverse weather is scarce. Synthetic
data offers a scalable, automatically labeled alternative, yet models trained on it
suffer a \emph{Sim2Real gap} when deployed on real imagery, especially under rain,
glare, fog/snow and night conditions.

This thesis studies and closes that gap for a YOLOv8-pose model that detects the four
runway corners, fine-tuned on nominal-condition synthetic data from LARD-V2 and
evaluated on a weather-stratified real test set from LARD-V1. It first quantifies the
gap through a controlled, shared-location cross-evaluation of five synthetic rendering
sources and a combined source, identifying which source generalizes best overall and
per weather scenario, and isolating a runway orientation/distance bias that
confounds naive real-versus-synthetic comparisons.

It then evaluates three complementary techniques, ordered from the least to the most
direct intervention on the domain shift: input-space weather augmentation with a
ratio study; self-training with confidence-filtered pseudo-labels on unlabeled real
landing sequences, extended with tracker-based temporal consistency; and
label-preserving, diffusion-based style transfer (Cosmos Transfer~2.5), optionally
compared against a CycleGAN baseline. All strategies are compared under one
lightweight, reproducible protocol (a fine-tuned YOLOv8-pose model, a fixed
5000-image budget, and a weather-stratified real test set), reporting primarily
mAP@[.50:.95] together with IoU-based localization and inference speed.

We find that the gap is only measurable once the test sets are matched on runway
orientation as well as size, that it concentrates in adverse weather (fog being the
hardest), and that self-training on unlabeled real footage is the most effective
technique so far --- improving adverse-weather splits by over ten mAP points --- while
weather augmentation helps only modestly and diffusion-based style transfer is ongoing.
